
Low-rank adaptation (LoRA) has become the standard method for parameter-efficient fine-tuning, yet practitioners universally default to rank $r=16$ without principled guidance. This work presents a systematic study of how optimal LoRA rank scales with model size. Across Pythia models (1B--12B parameters) on extractive QA tasks, optimal rank follows a sub-linear power law: $r_{\text{opt}} \propto N^{\alpha}$ with $\alpha = 0.82$ (95\% CI: [0.73, 0.92]) on SQuAD-v2. The relationship exhibits phase transition behavior: rank sensitivity at 12B is approximately $2\times$ higher than at 1B (ratio $= 2.08$, 95\% CI: [1.30, 3.74]). However, the scaling exponent is task-dependent---multi-hop reasoning (HotpotQA) yields $\alpha = 0.30$ (95\% CI: [0.09, 0.45]), substantially different from single-hop QA. Contrary to initial predictions, attention entropy decreases with model size (Pearson $r = -0.9999$, $p = 0.01$), suggesting larger models develop focused rather than distributed attention patterns. These results demonstrate that optimal LoRA rank scales sub-linearly with model size but that the specific relationship depends on task characteristics.
